% Partie : applied-modal-logics
% Chapitre : epistemic-logic
% Section : bisimulations

\documentclass[../../../include/open-logic-section]{subfiles}

\begin{document}

\olfileid[fr]{aml}{el}{bsd}

\olsection{Bisimulations}

Une question subsiste sur le pouvoir expressif du langage épistémique :
comment les modèles sont-ils liés aux formules qui y sont vraies ?
Les résultats de correspondance montrent que la validité de certaines
formules sur un cadre impose à celui-ci certaines propriétés. Mais,
à valuations propositionnelles identiques en tous les mondes, notre
langage modal permet-il de distinguer un modèle réduit à un monde
muni d'une boucle réflexive d'une chaîne infinie de mondes dont chacun
n'a que le suivant pour successeur ? Existe-t-il une formule $A$ qui
soit vraie à un seul de leurs deux mondes de départ ?

La bisimulation est une relation entre modèles relationnels qui
exprime un accord sur les valuations et sur les possibilités de
transition. Comme nous le verrons, elle donne une notion d'équivalence
qui préserve ce que notre langage épistémique peut exprimer ; elle
n'exige pas que les graphes des modèles soient identiques.

\begin{defn}[Bisimulation]
Soient $M_1 = \tuple{W_1, R_1, V_1}$ et $M_2 = \tuple{W_2, R_2, V_2}$
deux modèles relationnels pour le même ensemble d'agents, et\footnote{L'ensemble d'agents, noté A sans introduction dans ces clauses de la source, est ici noté G, conformément au langage du chapitre. Les comparaisons illustrées supposent l'accord des valuations aux mondes appariés.}
$\mathcal{R} \subseteq W_1 \times W_2$ une relation binaire.
On dit que $\mathcal{R}$ est une \emph{bisimulation} si, pour chaque
$\tuple{w_1, w_2} \in \mathcal{R}$, les conditions suivantes sont vérifiées :

\begin{enumerate}
  \item $w_1 \in V_1(p)$ si et seulement si $w_2 \in V_2(p)$,
  pour toute !!{propositional variable}~$p$.
  \item Pour tout agent $a \in G$ et tout monde $v_1 \in W_1$, si
  $R_{1_a} w_1 v_1$, il existe un monde $v_2 \in W_2$ tel que
  $R_{2_a} w_2 v_2$ et $\tuple{v_1, v_2} \in
  \mathcal{R}$.
  \item Pour tout agent $a \in G$ et tout monde $v_2 \in W_2$, si
   $R_{2_a} w_2 v_2$, il existe un monde $v_1 \in W_1$ tel que
   $R_{1_a} w_1 v_1$ et $\tuple{v_1, v_2} \in
   \mathcal{R}$.
\end{enumerate}

Lorsqu'une bisimulation entre $M_1$ et $M_2$ relie les mondes
$w_1$ et $w_2$, on écrit aussi $\tuple{M_1, w_1} \leftrightarroweq
\tuple{M_2, w_2}$ et on dit que $\tuple{M_1, w_1}$ et
$\tuple{M_2, w_2}$ sont \emph{bisimilaires}.
\end{defn}

Les clauses de la définition garantissent différentes propriétés.
La première garantit que deux mondes bisimilaires satisfont les mêmes
formules sans modalités, puisqu'ils s'accordent sur toutes les
!!{propositional variable}s. Les deux autres, souvent appelées
respectivement « forth » et « back » (aller et retour), garantissent
que les transitions accessibles pour chaque agent peuvent être
reproduites dans les deux sens entre mondes reliés.

\begin{thm}
Si $\tuple{M_1, w_1} \leftrightarroweq \tuple{M_2, w_2}$, alors,
pour toute !!{formula}~$!A$, on a $\mSat{M_1}{!A}[w_1]$ si et
seulement si $\mSat{M_2}{!A}[w_2]$.
\end{thm}

\begin{figure}
  \begin{center}
    \begin{tikzpicture}[modal]
      \node[world] (w1) {$w_1$};
      \node[world] (w2) [above left=of w1]{$w_2$};
      \node[world] (w3) [above right=of w1] {$w_3$};
      \draw[<->] (w1) to node {$a$} (w2);
      \draw[->,loop below] (w1) to node {$a$} (w1);
      \draw[<->] (w1) to [swap] node {$a$} (w3);
      \draw[->,loop above] (w2) to node {$a$} (w2) ;
      \draw[->,loop above] (w3) to node {$a$} (w3) ;

      \node[world](v1)[right of=w1, xshift=1.5in]{$v_1$};
      \node[world](v2)[above of=v1]{$v_2$};
      \draw[<->] (v1) to node {$a$} (v2);
      \draw[->,loop below] (v1) to node {$a$} (v1);
      \draw[->,loop above] (v2) to node {$a$} (v2) ;

      \draw[-,dotted] (w1) to (v1) ;
      \draw[-,dotted,bend left=45] (w2) to (v2) ;
      \draw[-,dotted,bend left=30] (w3) to (v2) ;
    \end{tikzpicture}
  \end{center}
  \caption{Deux modèles bisimilaires, si les valuations concordent
    aux mondes reliés par les pointillés.}
  \ollabel{fig:bisimilar}
\end{figure}

Les modèles de la \olref{fig:bisimilar} sont différents, mais, sous
la condition d'accord des valuations indiquée dans la légende,
une bisimulation relie les mondes $w_1$ et~$v_1$. Elle relie aussi
chacun des mondes $w_2$ et $w_3$ au monde~$v_2$ : rien dans notre
langage modal ne permet de distinguer ces mondes ainsi appariés.
La situation serait différente si $w_2$ et~$w_3$ satisfaisaient
des !!{propositional variable}s différentes.

\end{document}
